\justifying
\NSFCBodyText

按当前事实档案，申请人正在承担且与本课题方法或资料处理有关的项目如下。项目在研状态以 2026 年申报时点列示，正式提交前核对时点变化。

\textbf{（1）融合多源数据大模型驱动的光伏发电功率预测方法研究。}
自治区高校基本科研业务费科研项目—培育类，批准号 XJEDU2026P005，资助金额 5 万元，其中直接经费 3.9 万元、间接经费 1.1 万元，执行期 2026 年 1—12 月。申请人主持，负责总体方案、模型研究和成果组织。该项目以光伏功率为目标，本课题以 CIESM 气象场跨尺度重建为目标，可复用时序建模经验；不将气候场增强工作或拟产出的数据重复列为该项目已完成成果。

\textbf{（2）多模态融合技术研究及其在新疆产业集群中的应用。}
自治区天池英才引进计划（青年博士人才），依据任务书无独立项目批准号，资助金额 50 万元，执行期 2025 年 2 月至 2028 年 2 月。申请人主持，负责总体研究方案、关键模型设计和成果组织。其多源信息融合经验可支持本课题的方法实施；本课题另以 CIESM 数据、物理范围与聚合一致性作为研究任务和评价对象。

\textbf{（3）数据驱动下基于互补型生成对抗网络与神经统计模型的公路沙害智能预警研究。}
国家自然科学基金地区项目，批准号 62466056，总资助金额 32 万元，执行期 2025 年 1 月至 2028 年 12 月，基金委系统登记单位为新疆大学。申请人为参与者，负责数据收集整理。该经验可支持区域资料组织；本课题研究气候模式数据增强，不承担该项目公路沙害预警模型的重复研发任务。

本课题主要参与者的在研项目情况，将随\textbf{【待补 F-CIESM-02：参与者名单与分工】}一并核实，并按相关性补入正式申请材料。
